% Source-aligned English for printed p.174.
\SourcePage{174}

\LayerTarget{tohoku:application:leray-algebraic-varieties:p2}
\noindent of algebraic varieties (with their Zariski topology).

\LayerTarget{tohoku:remark:relative-leray-scope}
We shall confine ourselves here to this statement of the natural conditions
under which the Leray spectral sequence is valid, and shall not pursue its
study further.

\LayerTarget{tohoku:section:3.8}
\textbf{3.8. Comparison with Čech cohomology.} We refer to [15] for the
definition of the ``cohomology groups'' of $X$ with coefficients in an
abelian sheaf $F$, calculated by the method of Čech coverings.

\LayerTarget{tohoku:def:cech-cohomology}
We denote these groups by $\check H^p(X,F)$, to distinguish them from the
groups $H^p(X,F)$ defined in No.~2. (For simplicity, we shall consider no
``family $\Phi$'' other than the family of all closed subsets.) We note,
however, that these groups can be defined when $F$ is assumed only to be a
\textit{presheaf} of abelian groups:

\LayerTarget{tohoku:construction:cech-cochain-complex}
for every open covering $\mathcal U=(U_i)$ of $X$, one can form the complex

\LayerTarget{tohoku:eq:cech-cochain-complex-total}
\[
C(\mathcal U,F)=\sum_p C^p(\mathcal U,F)
\]
of cochains of $\mathcal U$ with values in the presheaf $F$, and set

\LayerTarget{tohoku:def:cover-cohomology}
\[
H^p(\mathcal U,F)=H^p(C(\mathcal U,F)),
\]
then take

\LayerTarget{tohoku:eq:cech-cohomology-direct-limit}
\[
\check H^p(X,F)=\varinjlim_{\mathcal U} H^p(\mathcal U,F),
\]
where the direct limit is taken over the filtered partially ordered set of
``classes of open coverings'' of $X$ (two open coverings being regarded as
equivalent if each refines the other).

\LayerTarget{tohoku:obstruction:cech-not-cohomological}
Unfortunately, the $\check H^p(X,F)$ do not in general form a cohomological
functor on the category $C^X$ of sheaves of abelian groups on $X$ (see the
example at the end of this section).

\LayerTarget{tohoku:construction:cech-delta-functor}
But $(\check H^0,\check H^1)$ forms an exact $\partial$-functor [15], [11].
Moreover, the $\check H^p$ are effaceable functors for $p>0$:

\LayerTarget{tohoku:lemma:product-sheaf-cech-acyclic}
for this, it suffices, for example, to show that if $M$ is the product sheaf
defined by a family $(M_x)_{x\in X}$ of abelian groups (cf. No.~1), then the
$\check H^p(X,M)$ vanish for $p>0$. In fact, one even proves
$H^p(\mathcal U,M)=0$ for every open covering $\mathcal U$ of $X$, using the
well-known homotopy operator employed classically when $M$ is fine and $X$ is
paracompact.

\LayerTarget{tohoku:comparison:cech-and-derived-cohomology}
\textit{It follows that, if $X$ is such that the $\check H^p$ can be regarded
as the components of a cohomological functor, then these $\check H^p$ are
canonically isomorphic to the functors $H^p$. This is so if $X$ is
paracompact} (cf., for example, [15, No.~25]), or if $X$ is arbitrary but one
restricts to the values $p=0$ and $p=1$ (as we already observed in No.~4).

\LayerTarget{tohoku:construction:cartan-leray-fixed-cover:p1}
More precise results are related to the Cartan--Leray spectral sequence of a
covering, as M. Cartan pointed out to me (whose idea I adopt here). Let
$\mathcal U$ be a fixed open covering of $X$, and put
$C(F)=C(\mathcal U,F)$ for every sheaf $F$. This gives a covariant left-exact
functor from $C=C^X$ to the category $C'$ of complexes of abelian groups in
positive degrees. Moreover, we have seen that the functors $H^p(C(F))$ are
effaceable. One readily concludes that there is a spectral functor abutting to
the right derived functor

\TohokuContinuation{175}
